\subsection{国内外研究现状}\label{sec:related-work}

机器人基础软件层面的研究可大致归为三个方向：操作系统、通信构件与开发 SDK；
本节按此三个方向梳理国内外现有工作，并指出研究空白。

在操作系统层面，研究核心是\textbf{实时性保证}。
标准 Linux 内核基于 CFS 公平调度，不提供硬实时保证；
PREEMPT-RT 补丁通过把内核大段代码改为可抢占，
将 ARM 嵌入式平台上的网络通信最坏延迟从 vanilla Linux 的 13197\,µs 压降至 88–110\,µs\cite{rtlinux}，
是目前机器人控制系统最主流的 Linux 实时化方案；
但仅有 PREEMPT-RT 还不够，并发流量下仍需 cpusets CPU 绑核隔离才能维持有界延迟\cite{rtlinux}。
ROS2 实时性研究综述\cite{ros2rtsurvey} 进一步指出，
内核 Executor 与定制调度器是该方向密集投入的研究分支。
近期 AMS\cite{ams}（Zheng et al., 2025）将 action context、action exception、action replay
三个 OS 原语融入 VLA 推理管理，在真实机械臂任务上将成功率提升 7×–24×，执行步数减少 29\%–74\%，
表明 OS 层机制对机器人推理的性能与可靠性具有直接影响。

在通信构件层面，ROS2\cite{ros2} 以 DDS（Data Distribution Service）为底层通信层，
是学术界与工业界的事实标准；但其节点调度仍依赖 Linux CFS，无硬实时保证。
针对 ARM 嵌入式平台的实测表明，
节点组合方式（intra-process vs.\ multi-process）对 CPU 占用与消息延迟有显著影响\cite{ros2composition}；
在树莓派3/4 等平台上，Linux 网络栈被发现是制约硬实时以太网通信的主要瓶颈\cite{ros2ddsrpi}。
为支撑此类分析，ros2\_tracing\cite{ros2tracing} 提供了基于 LTTng 的低开销追踪框架，
可在不显著扰动主循环的条件下获取 ROS2 节点调度与消息传递的精细时序。

在 SDK 层面，深度学习推理框架是机器人基础软件最核心的组成之一。
TFLite\cite{tflite} 面向 MCU 级极端受限场景，支持 INT8 量化与 ARM NEON SIMD 加速；
ONNX Runtime\cite{onnxruntime} 通过 KleidiAI 算子库在 ARM64 上实现 28–51\% 的性能提升；
PyTorch 则提供 TorchScript 与 CPU 优化，是 LeRobot\cite{lerobot} 等机器人开源框架的默认推理后端。
在机器人生态软件方面，HuggingFace LeRobot 等以 Python 为主的开源框架
快速降低了 ACT\cite{act}、Diffusion Policy\cite{diffusionpolicy} 等策略的训练与部署门槛，
但围绕其性能特性的研究仍以模型精度与吞吐量为主要评估维度。
更关键的是，上述推理框架与生态软件均以"外挂式"方式接入操作系统：
其内置线程池（OpenMP 或 C10）的调度与 OS 内核 CFS 调度器相互独立，
在资源受限平台上容易引发调度冲突，引入额外的 futex 阻塞与上下文切换开销，
而这一类内核层耦合开销迄今缺乏系统性量化。

综上，已有工作多聚焦于上述三个层面中的\textbf{单一层面}：
或针对内核做实时化改造、或评估通信中间件的延迟、或优化推理框架的算子性能；
即便涉及多层（如 ros2\_tracing 联合用户态—内核态追踪），其覆盖范围也仅限于 ROS2 通信链路。
\textbf{从 Python 运行时到 AI 推理框架再到 Linux 内核}这条贯穿机器人推理主循环的端到端调用链，
迄今仍缺乏系统性的跨层次联合剖析。
本文的研究方向正是\textbf{跨层次联合延迟分析}：
建立覆盖 AI 推理框架层、Python 运行时层、OS 内核层与硬件 I/O 层的四层性能模型，
通过ftrace等工具联合剖析，
端到端定位 ARM 嵌入式平台上机器人推理主循环的真实瓶颈所在。
